\documentclass{article}
\usepackage[UTF8]{ctex}
\usepackage{amsmath,amssymb,amsfonts}
\usepackage{enumitem}
\usepackage{geometry}
\geometry{a4paper, margin=1in}
% === 错题本排版优化（不参与题目标识） ===
\setlength{\parindent}{0pt}
\setlength{\parskip}{0.35em plus 0.10em minus 0.08em}
\setlength{\emergencystretch}{3em}
\setlist{itemsep=0.25em, topsep=0.40em, parsep=0pt}
\allowdisplaybreaks[2]
% === 错题本排版优化结束 ===
\begin{document}
\section*{陕西师范大学2024年数学分析真题}
\section*{187．陕西师范大学2024年数学分析试卷}

\subsection{题目1}
1．（15 分）判断对错并说明理由（正确的给出证明，错误的给出反例）．

\begin{itemize}
  \item[（1）] 若 \(f(x)\) 在 \((a, b)\) 上连续，则 \(f(x)\) 在 \((a, b)\) 上有界．
  \item[（2）] 若 \(f(x)\) 在 \((a, b)\) 上连续，且 \(\lim _{x \rightarrow a^{+}} f(x)\) 与 \(\lim _{x \rightarrow b^{-}} f(x)\) 均存在，则 \(f(x)\) 在 \((a, b)\) 上有界．
  \item[（3）] 若 \(f(x)\) 在 \((a, b)\) 上连续且有界，则 \(\lim _{x \rightarrow a^{+}} f(x)\) 与 \(\lim _{x \rightarrow b^{-}} f(x)\) 均存在．
\end{itemize}

\subsection{题目2}
2．（15 分）数列

\[
a_{1}=2, a_{2}=2+\frac{1}{2}, a_{3}=2+\frac{1}{2+\frac{1}{2}}, a_{4}=2+\frac{1}{2+\frac{1}{2+\frac{1}{2}}}, \cdots .
\]

问数列 \(\left\{a_{n}\right\}\) 是否收敛？说明理由．

\subsection{题目3}
3．（15 分）已知

\[
f(x)= \begin{cases}(x-1)^{3} \cos \frac{1}{x}, & x<1 \\ a, & x=1 \\ e^{b(x-1)}+c, & x>1\end{cases}
\]

问：\(a, b, c\) 为何值时，\(f(x)\) 在 \(x=1\) 处可导，并计算 \(f^{\prime}(1)\) 。

\subsection{题目4}
4．（15 分）已知 \(z=z(x, y)\) 是由方程 \(x+2 y+3 z=e^{z}\) 确定的隐函数，计算 \(\frac{\partial z}{\partial x}, \frac{\partial z}{\partial y}, \frac{\partial^{2} z}{\partial x^{2}}, \frac{\partial^{2} z}{\partial y \partial x}\) ．

\subsection{题目5}
5．（15 分）求抛物线 \(y^{2}=p x, y^{2}=q x(0<p<q)\) 及双曲线 \(x y=a, x y=b(0<a<b)\) 所围成区域 \(D\)的面积。

\subsection{题目6}
6．（15 分）计算

\[
\iint_{\sum} \frac{(x+y) \mathrm{d} y \mathrm{~d} z+y \mathrm{~d} z \mathrm{~d} x+2 z \mathrm{~d} x \mathrm{~d} y}{\sqrt{x^{2}+y^{2}+z^{2}}} .
\]

其中 \(\sum\) 为球面 \(x^{2}+y^{2}+z^{2}=a^{2}(a>0)\) 的外侧．

\subsection{题目7}
7．（15 分）设 \(f\) 在闭区间 \([a, b]\) 上有定义，且对任意的 \(x, y \in[a, b]\) ，有 \(|f(x)-f(y)| \leq 2(x-y)^{2}\) ，证明：\(f(x)\) 在 \([a, b]\) 上为常值函数．

\subsection{题目8}
8．（15 分）求幂级数 \(\sum_{n=1}^{\infty} \frac{\ln (n+1)}{n} x^{n}\) 的收敛域．

\subsection{题目9}
9．（15 分）设三元函数 \(f(x, y, z)\) 在 \(D=\{(x, y, z) \mid x+y+z \leq 1\}\) 上可微，且 \(\forall(x, y, z) \in D\) ，有

\[
\left|f_{x}(x, y, z)\right| \leq 1,\left|f_{y}(x, y, z)\right| \leq 2,\left|f_{z}(x, y, z)\right| \leq 3 .
\]

证明：\(\forall\left(x_{1}, y_{1}, z_{1}\right),\left(x_{2}, y_{2}, z_{2}\right) \in D\) ，有

\[
\left|f\left(x_{2}, y_{2}, z_{2}\right)-f\left(x_{1}, y_{1}, z_{1}\right)\right| \leq\left|x_{2}-x_{1}\right|+2\left|y_{2}-y_{1}\right|+3\left|z_{2}-z_{1}\right| .
\]

\subsection{题目10}
10．（15 分）设 \(f(x)\) 在闭区间 \([a, b]\) 上只有可去间断点，证明：\(f(x)\) 在 \([a, b]\) 上有界．
\end{document}
